%% TOAN12_C5_ONTAP1.tex -- Đề ôn tập chương V, đề số 1
%% Theo cấu trúc đề thi tốt nghiệp THPT: 12 TN, 4 Đ/S, 6 trả lời ngắn.
%% Đáp án: \dapan{khoá}{hướng dẫn}.
\begin{deontap}{1}
\phanI
\begin{cau}Trong không gian $Oxyz$, đường thẳng $\Delta\colon \begin{cases}x=-3-3t\\y=-3+t\\z=3-3t\end{cases}$ đi qua điểm nào dưới đây?
\pa{$M(3;3;-3)$}{$M(-3;-3;3)$}{$M(-3;1;-3)$}{$M(-3;-3;4)$}
\dapan{B}{Với $t=0$ ta được điểm $(-3;-3;3)$.}\end{cau}
\begin{cau}Trong không gian $Oxyz$, điểm nào dưới đây thuộc mặt phẳng $(P)\colon -x-3y+3z-5=0$?
\pa{$Q(1;3;5)$}{$Q(2;3;4)$}{$Q(2;4;4)$}{$Q(1;3;6)$}
\dapan{A}{Thay toạ độ vào phương trình $(P)$.}\end{cau}
\begin{cau}Trong không gian $Oxyz$, phương trình tham số của đường thẳng đi qua điểm $M(1;3;2)$ và có vectơ chỉ phương $\vec u=(1;-1;-3)$ là
\pa{$\begin{cases}x=1+t\\y=3-t\\z=2-3t\end{cases}$}{$\begin{cases}x=1+t\\y=-1+3t\\z=-3+2t\end{cases}$}{$\begin{cases}x=1\\y=3+2t\\z=2+4t\end{cases}$}{$\begin{cases}x=-1+t\\y=-3-t\\z=-2-3t\end{cases}$}
\dapan{A}{Áp dụng định nghĩa phương trình tham số.}\end{cau}
\begin{cau}Trong không gian $Oxyz$, trục $Ox$ có một vectơ chỉ phương là
\pa{$\vec u=(0;0;1)$}{$\vec u=(1;1;1)$}{$\vec u=(1;0;0)$}{$\vec u=(0;1;0)$}
\dapan{C}{}\end{cau}
\begin{cau}Trong không gian $Oxyz$, cho ba điểm $A(-2;1;2)$, $B(-1;-2;-2)$, $C(-1;-1;1)$. Phương trình mặt phẳng $(ABC)$ là
\pa{$5x+3y-z+10=0$}{$x-3y-4z+13=0$}{$5x+3y-z+9=0$}{$4x+3y-z+7=0$}
\dapan{C}{$\vt{AB}=(1;-3;-4)$, $\vt{AC}=(1;-2;-1)$, $\vec n=[\vt{AB},\vt{AC}]=(-5;-3;1)$; $(ABC)\colon 5x+3y-z+9=0$.}\end{cau}
\begin{cau}Trong không gian $Oxyz$, góc giữa đường thẳng $\Delta$ có vectơ chỉ phương $\vec u=(2;1;-2)$ và mặt phẳng $(P)\colon 2x+y-2z+1=0$ bằng
\pa{$90^\circ$}{$0^\circ$}{$60^\circ$}{$120^\circ$}
\dapan{A}{$\sin(\Delta,(P))=\dfrac{9}{9}$ nên góc bằng $90^\circ$.}\end{cau}
\begin{cau}Trong không gian $Oxyz$, góc giữa trục $Oz$ và mặt phẳng $(P)\colon x+2y-z+1=0$ (làm tròn đến hàng phần mười của độ) xấp xỉ bằng
\pa{$12^\circ$}{$24{,}1^\circ$}{$65{,}9^\circ$}{$39{,}1^\circ$}
\dapan{B}{$\sin=\dfrac{\sqrt{6}}{6}$.}\end{cau}
\begin{cau}Trong không gian $Oxyz$, mặt phẳng đi qua điểm $M(0;-2;0)$, song song với trục $Ox$ và vuông góc với mặt phẳng $(Q)\colon x+3y-z+5=0$ có phương trình là
\pa{$x+z=0$}{$y+3z+2=0$}{$x+3y-z+6=0$}{$y+3z-2=0$}
\dapan{B}{$\vec n=[\vec i,\vec n_Q]=(0;1;3)$; $y+3z+2=0$.}\end{cau}
\begin{cau}Trong không gian $Oxyz$ (đơn vị: km), một máy bay cất cánh từ điểm $A(-37;37;0)$ và bay theo đường thẳng có vectơ chỉ phương $\vec u=(4;2;1)$. Khi máy bay đạt độ cao $20$ km (tức $z=20$) thì nó ở vị trí $M(a;b;20)$. Giá trị $a+b$ bằng
\hinh{\begin{tikzpicture}[scale=0.85,font=\footnotesize,line join=round,line cap=round,>=Stealth]
\draw[->,gray] (0.000,0.000)--(-1.440,-1.080) node[below left]{$x$};
\draw[->,gray] (0.000,0.000)--(5.200,0.000) node[right]{$y$};
\draw[->,gray] (0.000,0.000)--(0.000,3.400) node[above]{$z$};
\node[below right,gray] at (0.000,0.000){$O$};
\fill[gray!10] (-0.300,0.000)--(-1.596,-0.972)--(3.504,-0.972)--(4.800,0.000)--cycle;
\draw[thick] (-0.336,0.298)--(3.768,2.476); \draw[->,thick] (-0.336,0.298)--(1.716,1.387);
\fill (-0.336,0.298) circle (.05) node[below right]{$A$}; \fill (3.768,2.476) circle (.05) node[above]{$B$};
\draw[dashed] (-0.336,0.298)--(-0.336,-0.702) (3.768,2.476)--(3.768,-0.324);

\end{tikzpicture}}

\pa{$120$}{$20$}{$-34$}{$86$}
\dapan{A}{Đường bay: $\begin{cases}x=-37+4t\\ y=37+2t\\ z=t\end{cases}$. $z=20\Leftrightarrow t=20$, $M(43;77;20)$, $a+b=120$.}\end{cau}
\begin{cau}Trong không gian $Oxyz$ (đơn vị: mét), mặt một sườn đồi được mô hình hoá bởi mặt phẳng $(P)\colon 4x+8y+z-7=0$. Một flycam đang ở vị trí $F(1;0;12)$. Khoảng cách (mét) từ flycam đến mặt sườn đồi là
\hinh{\begin{tikzpicture}[scale=0.85,font=\footnotesize,line join=round,line cap=round,>=Stealth]
\draw[->,gray] (0.000,0.000)--(-1.440,-1.080) node[below left]{$x$};
\draw[->,gray] (0.000,0.000)--(5.000,0.000) node[right]{$y$};
\draw[->,gray] (0.000,0.000)--(0.000,3.300) node[above]{$z$};
\node[below right,gray] at (0.000,0.000){$O$};
\fill[gray!10] (-0.500,0.000)--(-1.760,-0.945)--(3.240,-0.945)--(4.500,0.000)--cycle;
\draw[dashed] (0.300,1.525)--(0.300,-0.675); \fill (0.300,-0.675) circle (.04);
\filldraw[fill=gray!40] (0.300,1.525) ++(-.18,0) rectangle ++(.36,.1); \node[above] at (0.300,1.525){$A$};
\draw[dashed] (3.240,2.530)--(3.240,-0.270); \fill (3.240,-0.270) circle (.04);
\filldraw[fill=gray!40] (3.240,2.530) ++(-.18,0) rectangle ++(.36,.1); \node[above] at (3.240,2.530){$B$};
\draw[thick,dotted] (0.300,1.525)--(3.240,2.530);
\end{tikzpicture}}

\pa{$1$ m}{$9$ m}{$\dfrac{23}{9}$ m}{$2$ m}
\dapan{A}{$d(F,(P))=\dfrac{|9|}{9}=1$ (m).}\end{cau}
\begin{cau}Trong một hệ trục $Oxyz$ (đơn vị: mét), một tấm kính phẳng được đặt đi qua ba điểm $A(7;0;0)$, $B(0;4;0)$, $C(0;0;4)$. Một đèn nhỏ đặt tại điểm $E(4;1;2)$. Khoảng cách (mét) từ đèn đến tấm kính (làm tròn đến hàng phần trăm) là
\hinh{\begin{tikzpicture}[scale=0.85,font=\footnotesize,line join=round,line cap=round,>=Stealth]
\draw[->,gray] (0.000,0.000)--(-1.260,-0.945) node[below left]{$x$};
\draw[->,gray] (0.000,0.000)--(4.300,0.000) node[right]{$y$};
\draw[->,gray] (0.000,0.000)--(0.000,3.200) node[above]{$z$};
\node[below right,gray] at (0.000,0.000){$O$};
\fill[gray!25,opacity=.8] (-0.936,-0.702)--(3.400,0.000)--(0.000,2.500)--cycle; \draw[thick] (-0.936,-0.702) node[left]{$A$}--(3.400,0.000) node[below]{$B$}--(0.000,2.500) node[left]{$C$}--cycle;
\shade[ball color=gray!20] (0.640,0.030) circle (.1); \node[right] at (0.640,0.030){$E$};
\end{tikzpicture}}

\pa{$0{,}84$ m}{$0{,}2$ m}{$1{,}69$ m}{$1{,}34$ m}
\dapan{A}{Mặt phẳng $(ABC)\colon \dfrac{x}{7}+\dfrac{y}{4}+\dfrac{z}{4}=1\Leftrightarrow 4x+7y+7z-28=0$. $d=\dfrac{|36|}{\sqrt{1824}}\approx 0{,}84$.}\end{cau}
\begin{cau}Trong không gian $Oxyz$ (đơn vị: km), một máy bay cất cánh từ điểm $A(-46;-4;0)$ và bay theo đường thẳng có vectơ chỉ phương $\vec u=(4;-4;2)$. Khi máy bay đạt độ cao $10$ km (tức $z=10$) thì nó ở vị trí $M(a;b;10)$. Giá trị $a+b$ bằng
\hinh{\begin{tikzpicture}[scale=0.85,font=\footnotesize,line join=round,line cap=round,>=Stealth]
\draw[->,gray] (0.000,0.000)--(-1.440,-1.080) node[below left]{$x$};
\draw[->,gray] (0.000,0.000)--(5.200,0.000) node[right]{$y$};
\draw[->,gray] (0.000,0.000)--(0.000,3.400) node[above]{$z$};
\node[below right,gray] at (0.000,0.000){$O$};
\fill[gray!10] (-0.300,0.000)--(-1.596,-0.972)--(3.504,-0.972)--(4.800,0.000)--cycle;
\draw[thick] (-0.336,0.298)--(3.768,2.476); \draw[->,thick] (-0.336,0.298)--(1.716,1.387);
\fill (-0.336,0.298) circle (.05) node[below right]{$A$}; \fill (3.768,2.476) circle (.05) node[above]{$B$};
\draw[dashed] (-0.336,0.298)--(-0.336,-0.702) (3.768,2.476)--(3.768,-0.324);

\end{tikzpicture}}

\pa{$-50$}{$-45$}{$-2$}{$-52$}
\dapan{A}{Đường bay: $\begin{cases}x=-46+4t\\ y=-4-4t\\ z=2t\end{cases}$. $z=10\Leftrightarrow t=5$, $M(-26;-24;10)$, $a+b=-50$.}\end{cau}
\phanII
\begin{cau}Trong không gian $Oxyz$, cho ba điểm $A(1;0;0)$, $B(0;2;0)$, $C(0;0;2)$.
\begin{dungsai}
  \y Mặt phẳng $(ABC)$ có phương trình $x+2y+2z=1$. \oDS
  \y Vectơ $\vec n=(1;2;2)$ là một vectơ pháp tuyến của $(ABC)$. \oDS
  \y Khoảng cách từ $O$ đến $(ABC)$ bằng $\dfrac{\sqrt{6}}{3}$. \oDS
  \y Thể tích khối tứ diện $OABC$ bằng $\dfrac{2}{3}$. \oDS
\end{dungsai}
\dapan{S,S,Đ,Đ}{$(ABC)\colon 2x+y+z-2=0$; $\mathrm d(O,(ABC))=\dfrac{\sqrt{6}}{3}$; $V=\dfrac{1\cdot2\cdot2}6$.}\end{cau}
\begin{cau}Trong không gian $Oxyz$, cho mặt phẳng $(P)\colon x+z-3=0$.
\begin{dungsai}
  \y Một vectơ pháp tuyến của $(P)$ là $(0;1;-3)$. \oDS
  \y Góc giữa đường thẳng $\Delta$ có vectơ chỉ phương $(0;0;1)$ và $(P)$ bằng $45^\circ$. \oDS
  \y Góc giữa $(P)$ và mặt phẳng $(Q)\colon x+y=0$ bằng $30^\circ$. \oDS
  \y Góc giữa $(P)$ và mặt phẳng $(R)\colon -x+z+5=0$ bằng $90^\circ$. \oDS
\end{dungsai}
\dapan{S,Đ,S,Đ}{$\sin(\Delta,(P))=|\cos(\vec u,\vec n)|$; $\cos((P),(Q))=|\cos(\vec n_P,\vec n_Q)|$.}\end{cau}
\begin{cau}Trong không gian $Oxyz$, cho hai mặt phẳng $(P)\colon 2x-y+2z+3=0$ và $(Q)\colon 2x-y+2z-6=0$.
\begin{dungsai}
  \y Hai mặt phẳng $(P)$ và $(Q)$ cắt nhau. \oDS
  \y Khoảng cách giữa $(P)$ và $(Q)$ bằng $3$. \oDS
  \y Khoảng cách từ $M(1;1;1)$ đến $(P)$ bằng $3$. \oDS
  \y Mặt phẳng $(R)\colon -x-2y+1=0$ vuông góc với $(P)$. \oDS
\end{dungsai}
\dapan{S,Đ,S,Đ}{Cùng vectơ pháp tuyến $(2;-1;2)$, $D$ khác nhau; $d=\dfrac{|3-(-6)|}{3}$; $(R)$ có $\vec n_R\cdot\vec n_P=0$.}\end{cau}
\begin{cau}Một đường dốc cho xe lăn được mô hình hoá trong hệ trục $Oxyz$ (đơn vị: m): chân dốc $A(0;0;0)$, đỉnh dốc $B(24;0;2)$, mặt đất là $(Oxy)$. Xét tính đúng sai của các mệnh đề sau:
\hinh{\begin{tikzpicture}[scale=0.85,font=\footnotesize,line join=round,line cap=round,>=Stealth]
\fill[gray!25] (0,0)--(5,0)--(5,.9)--cycle; \draw[thick] (0,0)--(5,.9)--(5,0)--cycle;
\node[below left] at (0,0){$A$}; \node[above right] at (5,.9){$B$}; \draw (.9,0) arc (0:10.2:.9); \node at (1.25,.12){$\alpha$};
\draw[<->] (5.25,0)--(5.25,.9) node[midway,right]{độ cao}; \draw[<->] (0,-.3)--(5,-.3) node[midway,below]{chiều dài theo phương ngang};
\end{tikzpicture}}

\begin{dungsai}
  \y Độ dài mặt dốc $AB$ bằng $2 \sqrt{145}$ m. \oDS
  \y Sin của góc giữa mặt dốc và mặt đất bằng $\dfrac{\sqrt{145}}{145}$. \oDS
  \y Theo quy chuẩn, đường dốc cho xe lăn có độ dốc (tỉ số chiều cao trên chiều dài nằm ngang) không vượt quá $\dfrac{1}{12}$. Đường dốc này đạt quy chuẩn. \oDS
  \y Góc giữa mặt dốc và mặt đất xấp xỉ $85{,}2^\circ$. \oDS
\end{dungsai}
\dapan{Đ,Đ,Đ,S}{$\overrightarrow{AB}=(24;0;2)$, mặt đất có vectơ pháp tuyến $(0;0;1)$, $\sin\alpha=\dfrac{2}{\sqrt{580}}$. a) Đúng. b) Đúng. c) Đúng. d) Sai: đó là phần bù của góc cần tìm.}\end{cau}
\phanIII
\begin{cau}Trong không gian $Oxyz$, tính bán kính của mặt cầu $(S)\colon x^2+y^2+z^2+4x+6y-6z+13=0$.\otraloi
\dapan{3}{$R=\sqrt{22-(13)}=3$.}\end{cau}
\begin{cau}Trong một căn phòng, chọn hệ trục $Oxyz$ (đơn vị mét) sao cho sàn nhà thuộc mặt phẳng $(Oxy)$ và một bức tường nằm trong mặt phẳng $(P)\colon 3x+4y-24=0$. Một chiếc đèn được treo tại điểm $D(1;2;4)$. Tính khoảng cách từ đèn đến bức tường (đơn vị mét).\otraloi
\dapan{2{,}6}{$d=\dfrac{|3+8-24|}{5}=2{,}6$.}\end{cau}
\begin{cau}Trong không gian $Oxyz$, gọi $M'(a;b;c)$ là điểm đối xứng của điểm $M(1;-1;2)$ qua mặt phẳng $(P)\colon x+2y+2z-12=0$. Tính $a+b+c$.\otraloi
\dapan{12}{Hình chiếu $H(2;1;4)$ là trung điểm $MM'$; $M'(3;3;6)$.}\end{cau}
\begin{cau}Trong không gian $Oxyz$, gọi $M'(a;b;c)$ là điểm đối xứng của điểm $M(1;3;-2)$ qua mặt phẳng $(P)\colon x+2y+2z+6=0$. Tính $a+b+c$.\otraloi
\dapan{-8}{Hình chiếu $H(0;1;-4)$ là trung điểm $MM'$; $M'(-1;-1;-6)$.}\end{cau}
\begin{cau}Trong một hệ trục $Oxyz$ (đơn vị: mét), một tấm kính phẳng được đặt đi qua ba điểm $A(8;0;0)$, $B(0;5;0)$, $C(0;0;4)$. Một đèn nhỏ đặt tại điểm $E(6;4;1)$. Khoảng cách (mét) từ đèn đến tấm kính bằng bao nhiêu (làm tròn kết quả đến hàng phần trăm)?
\hinh{\begin{tikzpicture}[scale=0.85,font=\footnotesize,line join=round,line cap=round,>=Stealth]
\draw[->,gray] (0.000,0.000)--(-1.260,-0.945) node[below left]{$x$};
\draw[->,gray] (0.000,0.000)--(4.300,0.000) node[right]{$y$};
\draw[->,gray] (0.000,0.000)--(0.000,3.200) node[above]{$z$};
\node[below right,gray] at (0.000,0.000){$O$};
\fill[gray!25,opacity=.8] (-0.936,-0.702)--(3.400,0.000)--(0.000,2.500)--cycle; \draw[thick] (-0.936,-0.702) node[left]{$A$}--(3.400,0.000) node[below]{$B$}--(0.000,2.500) node[left]{$C$}--cycle;
\shade[ball color=gray!20] (0.640,0.030) circle (.1); \node[right] at (0.640,0.030){$E$};
\end{tikzpicture}}
\otraloi
\dapan{2{,}33}{Mặt phẳng $(ABC)\colon \dfrac{x}{8}+\dfrac{y}{5}+\dfrac{z}{4}=1\Leftrightarrow 5x+8y+10z-40=0$. $d=\dfrac{|128|}{\sqrt{3024}}\approx 2{,}33$.}\end{cau}
\begin{cau}Một tấm pin mặt trời nằm trên mặt phẳng $(P)\colon 2x+2y+4z-4=0$ trong hệ trục $Oxyz$, mặt đất là mặt phẳng $(Oxy)$. Góc nghiêng của tấm pin so với mặt đất bằng bao nhiêu độ (làm tròn đến hàng đơn vị)?
\hinh{\begin{tikzpicture}[scale=0.9,font=\footnotesize,line join=round,line cap=round,>=Stealth]
\draw[thick] (-.3,0)--(4.3,0); \fill[pattern=north east lines] (-.3,-.15) rectangle (4.3,0);
\filldraw[fill=gray!50] (.6,0)--(3.2,1.5)--(3.32,1.32)--(.72,-.18)--cycle; \draw (2.2,0)--(2.2,.92); \draw (1.4,0) arc (0:30:.8); \node at (1.65,.2){$\varphi$};
\foreach \d in {0,.5,1}{\draw[->,gray] (1.6+\d,3.2)--(2.4+\d,2.4);} \node[gray,above] at (2.4,3.2){tia nắng};
\end{tikzpicture}}
\otraloi
\dapan{35}{$\cos\varphi=\dfrac{|4|}{\sqrt{24}}$, $\varphi\approx 35{,}3^\circ$.}\end{cau}
\end{deontap}
